% ===================================================================
% SKILL: UST HIZALI CIFT SUTUN, ANA GOVDE KUTUSUZ
% ===================================================================
% AMAC:
% Iki sutunlu yerlesimde ana anlatimi kutusuz kurmak, gerekiyorsa yalnizca
% kisa bir yardimci not kutusu kullanmak.
% ===================================================================

\begin{frame}[t]{Iki Sutunda Kutu Yerine Icerik}
\small
\begin{columns}[T,onlytextwidth]
    \begin{column}{0.48\textwidth}
        \textbf{Ayrik yapilar}

        \vspace{0.2cm}
        Sonuc kumesi sonlu veya sayilabilir sonsuz ise degisken ayriktir.

        \begin{itemize}
            \item zar sonucu
            \item hatali urun sayisi
            \item gelen cagri adedi
        \end{itemize}
    \end{column}
    \begin{column}{0.48\textwidth}
        \textbf{Surekli yapilar}

        \vspace{0.2cm}
        Olasi degerler bir araliktaki tum noktalari kapsiyorsa degisken sureklidir.

        \begin{itemize}
            \item bekleme suresi
            \item olcum hatasi
            \item sicaklik
        \end{itemize}
    \end{column}
\end{columns}

\vspace{0.25cm}
\begin{infoblock}{Kritik Ayirim}
    Sonsuz sayida deger tek basina surekli olmak anlamina gelmez; sayilabilir sonsuzluk yine ayrik olabilir.
\end{infoblock}
\end{frame}

\begin{frame}[t]{Karsilastirma ve Sonuc}
\small
\begin{columns}[T,onlytextwidth]
    \begin{column}{0.58\textwidth}
        \textbf{Okuma prensibi}

        \vspace{0.2cm}
        Ayrik yapida sayma mantigi, surekli yapida aralik ve alan mantigi baskindir.

        \[
            \text{ayrik} \Rightarrow \sum
            \qquad
            \text{surekli} \Rightarrow \int
        \]
    \end{column}
    \begin{column}{0.38\textwidth}
        \begin{resultblock}{Sonuc}
            Hangi dagilim aracini kullanacagin sorunun yapisi belirler.
        \end{resultblock}
    \end{column}
\end{columns}
\end{frame}
